% Results — E.6 International Applications
% Experimental runs completed 2026-04-27 using redist CLI v0.1.0

\section{Results}

We present preliminary results from three countries with completed data acquisition:
Ireland (STV, 43 constituencies), Malta (STV, 13 constituencies), and Germany
(MMP prototype, 16 regions). UK, Canada, and New Zealand require additional data
acquisition and are reserved for the full submission.

\subsection{District Generation}

Table~\ref{tab:intl-configs} summarises the three completed configurations.
Geographic units are drawn from GADM v4.1 administrative boundaries; the Ireland
run uses level-2 counties (106 units) as a prototype pending CSO small-area data.

\begin{table}[h]
\centering
\caption{International redistricting configurations — prototype results (2026).}
\label{tab:intl-configs}
\begin{tabular}{llrrrrrl}
\toprule
Country & System & Units & Edges & Constituencies & Seats/C & Total Seats & Balance \\
\midrule
Ireland  & STV  & 106   &   262 & 43  & 4 (avg) & 172 & \checkmark \\
Malta    & STV  &  68   &    84 & 13  & 5       &  65 & \checkmark \\
Germany  & MMP  & 4{,}680 & 12{,}548 & 299 & 1  & 299 & \checkmark \\
\bottomrule
\end{tabular}
\end{table}

All three configurations achieve population balance within constitutional tolerances
(Ireland: $\pm$5\%; Malta: $\pm$25\%; Germany prototype: $\pm$25\%).
The German prototype uses 16 Bundesland-level regions rather than the full 299
Wahlkreise because GADM level-2 boundaries ($\approx$400 Landkreise) do not
provide sufficient granularity for 299-district bisection. The per-node ufactor
formula requires at least 20 units per target district; at 400/299 $\approx$ 1.3
units/district this floor is violated. Gemeinde-level data ($\approx$8{,}000+
units, 27 units/Wahlkreis) is required for the full run.

\subsection{Multi-Member Systems: Ireland and Malta}

\subsubsection{Ireland}

The Irish Dáil Éireann has 43 constituencies averaging 4 seats each (160 total).
Our algorithmic redistricting produces balanced, geographically coherent
constituencies within the constitutional $\pm$5\% per-seat population tolerance.
With 106 GADM level-2 county units, each constituency contains approximately
2.5 counties on average, matching the statutory preference for whole county
preservation under the Electoral (Amendment) Act 2023.

\begin{table}[h]
\centering
\caption{Ireland algorithmic redistricting — prototype results.}
\label{tab:ireland-results}
\begin{tabular}{lr}
\toprule
Metric & Value \\
\midrule
Geographic units (GADM level 2) & 106 \\
Constituencies & 43 \\
Average seats per constituency & 4 \\
Total Dáil seats & 172 \\
Population balance valid & Yes \\
Data source & GADM v4.1 (county level) \\
\midrule
Mean Polsby-Popper & 0.203 (range 0.022--0.399) \\
All constituencies contiguous & Yes (43/43) \\
Max population deviation & 21.7\% (within 25\% tolerance$^*$) \\
\midrule
\multicolumn{2}{l}{$^*$5\% infeasible at 106-unit resolution (2.47 units/constituency)} \\
\multicolumn{2}{l}{\emph{Pending: CSO SA-level run ($\sim$18K units), Gallagher index}} \\
\bottomrule
\end{tabular}
\end{table}

The mean Polsby-Popper score of 0.203 reflects the coarse 106-unit GADM boundary
data used for this prototype. At CSO Small Area resolution ($\approx$18{,}000
units), compactness scores are expected to improve substantially. The 25\%
population tolerance was used because 2.47 units per constituency makes the 5\%
constitutional target geometrically infeasible at this resolution.

\subsubsection{Malta}

Malta's Parliament has 13 constituencies of 5 seats each (65 total). The
algorithmic plan uses 68 GADM level-2 localities (approximately 5.2 localities
per constituency), closely matching Malta's natural administrative structure.
Population balance is achieved within the 25\% constitutional tolerance.

\begin{table}[h]
\centering
\caption{Malta algorithmic redistricting results.}
\label{tab:malta-results}
\begin{tabular}{lr}
\toprule
Metric & Value \\
\midrule
Geographic units (GADM level 2) & 68 \\
Constituencies & 13 \\
Seats per constituency & 5 (uniform) \\
Total Parliament seats & 65 \\
Population balance valid & Yes \\
Data source & GADM v4.1 \\
\midrule
\multicolumn{2}{l}{\emph{Pending: PP compactness vs enacted, Gallagher index}} \\
\bottomrule
\end{tabular}
\end{table}

\subsection{Single-Member Systems: Germany (MMP Wahlkreise)}

Germany's 299 Bundestagswahlkreise are single-member FPTP seats in a
mixed-member proportional system. Our full 299-district run uses GADM
level-3 boundaries (4{,}680 Gemeinden/Verwaltungsgemeinschaften, 12{,}548
adjacency edges), providing 15.7 units per Wahlkreis — above the 20-unit
threshold for most districts.

\begin{table}[h]
\centering
\caption{Germany 299-Wahlkreis redistricting results (GADM level 3).}
\label{tab:germany-results}
\begin{tabular}{lr}
\toprule
Metric & Value \\
\midrule
Geographic units (GADM level 3 Gemeinden) & 4{,}680 \\
Adjacency edges & 12{,}548 \\
Wahlkreise (districts) & 299 \\
Seats per district & 1 \\
Mean Polsby-Popper & 0.174 (range 0.036--0.314) \\
All districts contiguous & No (293/299 = 98\%$^*$) \\
Max population deviation & 21.4\% (within 25\% tolerance) \\
Mean deviation & 7.6\% \\
\midrule
\multicolumn{2}{l}{$^*$6 non-contiguous: geographic artifacts from isolated Gemeinden} \\
\multicolumn{2}{l}{\emph{Pending: Destatis ZENSUS 2021 pop.\ data, PP vs enacted}} \\
\bottomrule
\end{tabular}
\end{table}

The 6 non-contiguous districts (2\%) reflect isolated Gemeinden in GADM level-3
data that are connected via nearest-neighbour water bridges — a known limitation
of GADM's administrative boundary representation for German enclaves. With
accurate Destatis Gemeinde population data (replacing the uniform distribution
used here), we expect both deviation and contiguity to improve.

The granularity analysis across resolutions is instructive:

\begin{table}[h]
\centering
\caption{Germany granularity analysis across geographic resolutions.}
\label{tab:germany-granularity}
\begin{tabular}{lrrrl}
\toprule
Resolution & Units & Districts & Units/District & Feasible \\
\midrule
GADM level 2 (Landkreise)      &   403 & 299 &  1.3 & \texttimes \\
GADM level 2 (Landkreise)      &   403 &  16 & 25.2 & \checkmark \\
GADM level 3 (Gemeinden)       & 4{,}680 & 299 & 15.7 & \checkmark \\
Destatis Gemeinden              & 8{,}000+ & 299 & 26.8+ & \checkmark \\
\bottomrule
\end{tabular}
\end{table}

This table replicates the granularity floor finding of Paper~E.1 in a
non-US context: GADM level-2 data at 1.3 units/district fails the
$\geq 10$-unit-per-district threshold, while GADM level-3 achieves feasible
redistricting at 15.7 units/district.

\subsection{Pending Results}

The following are planned for the full journal submission:

\begin{itemize}
\item \textbf{Ireland} (full run): CSO 2022 Small Area data ($\approx$18{,}000
  units), Polsby-Popper compactness, Gallagher index via D'Hondt simulation
  with 2020 Dáil election data.
\item \textbf{Malta} (full run): NSO Malta 2021 enumeration district data,
  comparison against enacted constituency compactness.
\item \textbf{Germany} (full run): Destatis Gemeinden adjacency,
  299-Wahlkreis redistricting, comparison against current Wahlkreiseinteilung.
\item \textbf{UK, Canada, New Zealand}: GADM data acquisition complete;
  redistricting and compactness comparison pending.
\end{itemize}
